%% 03-framework.tex — Conceptual Framework
%% Status: WRITTEN (v1) — concepts stable; expect light revision after results.
%% Source: NL_Robot_Synthesis_Short Part 1; Summary - Graduated Autonomy

\section{Method}
\label{sec:method}
\subsection{Desktop Robot}
We created a differential-drive desktop robot that runs ROS 2 Jazzy on a Raspberry Pi 5 with Dynamixel XL330 motors. It adapts a microphone (ReSpeaker Mic Array v2.0), and a plug-in speaker as the output. These are the only user-level communication channel. We intentional keep the feature of the robot minimal which can only adapt vocie input/output to comply with our study design purpose. Meanwhile, the embodiment is intentionally designed as round, which means we provide ambiguous toward the orientation, so participants should be able to communicate and negotiate with the robot.
(instruction: put desktop robot picture here)
\subsection{Deployment and Setting}
We recruited 12 participants from our campus facility. We conducted a disclosed wizard of Oz experiment with them, where they were told the natrue of the study when they sign the consent form. We request them to communicate with the robot directly, providing only the low level instruction. When they try to provide high level instruction such as "go toward me", or "go toward the vessel", the robot will provide. The finish task one followed by task two.
\subsubsection{Task one: exploration}
Task one composites a everday life desktop setting: laptop, vassel, cup, and other office tools were scattered on the desktop. The desktop robot sits at the other corner opposite of the participants, and the participants will need to use verbal instruction to make the robot to move toward them and meanwhile circumstance all the obstacles along the way.  
(instruction: Put study1 setting, and study one participants picture here)
\subsubsection{Task two: projected trajectory}
Task two requests participant to order the robot to move along the pre-defined trajectory and move toward themselves. Instead of the open-ended first task where participants have some improvisation space, this task require a stricter, concise level of wording description for participant to navigate the robot along the line. To give participant time to get familiar themselves with the task, the task is therefore designed in ascending order of difficulty. The first task is a simple right-angle turn, the second is the zig-zag line, and the final and thrid one is a complex route. 
(instruction: Put the study2 setting1-3 as pictures in here. Together)
\subsection{Video analysis}
Three researchers independly watched the video and come out with the coding scheme of the strategies that participants used for the interaction with the desktop robot. Then we gather together to finalized the coding that we are going to share. We created a browser-based video annotator, and each researcher independly coded these video, and then we output these file as 
(instruction: Put web interface here)

